% Source: https://artofproblemsolving.com/community/c3h1445027_fractions_with_a_denominator_of_7
% Stable identity: aops-topic-1445027-q01
For a number $x$ write $\{x\}$ for its fractional part. Multiplying a decimal by $10^j$ moves the decimal point $j$ places to the right, so the decimal digits of $\left\{10^j\cdot\dfrac{1}{7}\right\}$ are the digits of $\dfrac{1}{7}$ with the first $j$ digits removed; since the decimal of $\dfrac{1}{7}$ repeats with period $6$, that is the block $142857$ cyclically shifted by $j$ places.

Now $10^j=7q+r$ where $r$ is the remainder of $10^j$ mod $7$, so
\[
10^j\cdot\dfrac{1}{7}=q+\dfrac{r}{7},\qquad\text{and}\qquad \left\{10^j\cdot\dfrac{1}{7}\right\}=\dfrac{r}{7}.
\]
The remainders of $10^0,10^1,\dots,10^5$ mod $7$ are
\[
1,\ 3,\ 2,\ 6,\ 4,\ 5,
\]
which are all six of $1,2,\dots,6$. So every $\dfrac{k}{7}$ ($1\le k\le6$) equals $\left\{10^j\cdot\dfrac{1}{7}\right\}$ for some $j$ between $0$ and $5$, and its decimal is $0.\overline{142857}$ shifted by $j$ places. For example $10^2\equiv2$, and indeed $\dfrac{2}{7}=0.\overline{285714}$ is $142857$ shifted by $2$; $10^3\equiv6$ gives $\dfrac{6}{7}=0.\overline{857142}$. $\blacksquare$

The key fact is that $10$ has order $6$ modulo $7$, so its powers run through all nonzero remainders; the same argument works for every prime $p$ for which $10$ has order $p-1$ (for example $17$ and $19$).
